\section{Reading contract and overall map}
\label{sec:reading-contract}

\subsection{What this guide assumes and explains}

This guide assumes familiarity with rings, quotient rings, free modules,
equivalence relations and quotient sets, commutative diagrams, and induction.
It uses standard English mathematical terminology in introductions, headings,
figures, and the glossary, with consistent terms in repeated contexts.
It does not assume knowledge of the following.

\begin{itemize}
  \item The design of HAETAE or module-lattice signatures;
  \item Jasmin syntax and memory models;
  \item EasyCrypt's Hoare logic, pRHL, or tactic language;
  \item Implementation conventions for rANS entropy coding.
\end{itemize}

The goal is not a line-by-line explanation of every source file.
For each proof obligation, the guide asks: “What is the object?”,
“Which map must commute?”, “Which premises restrict its domain?”,
and “Which arrows have currently been proved for actual procedures?”
The project's scope is limited to mode~2 and pinned specification/implementation
snapshots \cite{project-scope}.

\begin{translation}[Translation across three layers]
Read the repository's theorems as propositions moving between three layers.
\begin{enumerate}
  \item \textbf{Algebraic layer}: \(R_q\)-module elements, decompositions,
        challenges, and paper algorithms;
  \item \textbf{Representation layer}: canonical byte strings, prefix
        projections, memory regions, and fixed-width integers;
  \item \textbf{Execution layer}: actual generated procedures, branches,
        loops, return values, and trace observations.
\end{enumerate}
Transporting an equality from one layer to another always requires a named
representation/refinement theorem. A shared mathematical intention does not
make a diagram commute automatically.
\end{translation}

\begin{figure}[ht]
\centering
\begin{tikzpicture}[node distance=8mm and 18mm]
  \node[layernode] (alg) {Paper algebraic layer\\\(R_q\)-modules\\and distributions};
  \node[layernode, right=of alg] (repr) {Canonical representation\\Byte strings\\and prefixes};
  \node[layernode, right=of repr] (proc) {Actual execution layer\\Jasmin\\state machines};
  \node[partialnode, below=12mm of repr] (security) {Top-level correctness\\and security\\composition goals};
  \draw[flowarrow] (alg) -- node[above,align=center,font=\small]{encode/\\refine} (repr);
  \draw[flowarrow] (repr) -- node[above,align=center,font=\small]{extract/\\simulate} (proc);
  \draw[openarrow] (proc) -- (security);
  \draw[openarrow] (alg) -- (security);
\end{tikzpicture}
\caption{Overall proof structure. Only some horizontal arrows are closed;
the top-level composition arrows remain partial.}
\label{fig:three-layers}
\end{figure}

\subsection{Status terms include logical quantifiers}

The operational source of truth is \sourcepath{CLAIM_LEDGER.md}
\cite{claim-ledger}; dependencies are recorded in
\sourcepath{THEOREM_GRAPH.md} \cite{theorem-graph}.
This guide uses status terms only with the following meanings.

\begin{center}
\begin{tabularx}{\textwidth}{L{28mm} Y Y}
\toprule
Status & Exact meaning & What does not follow automatically \\
\midrule
\statusproved & A named EasyCrypt declaration freshly compiles under the stated premises. &
Parent claims, termination, losslessness, distributional equality \\
\statuspartial & Child propositions or restricted directions are proved, but the named parent remains open. &
A global statement with the restrictions removed \\
\statusspecified & Operations, predicates, or interfaces expressing the goal compile. &
A proof of the goal \\
\statusblocked & A specific extraction, semantic, numerical, or probabilistic theorem is missing, preventing composition. &
A claim that the obstacle is insurmountable \\
\statusdeferred & Explicitly frozen at a paper boundary after a time limit. &
A claim that the obligation is false or unnecessary \\
\bottomrule
\end{tabularx}
\end{center}

In particular, “Hoare partial correctness is proved” means
\[
  \text{if the procedure returns, the postcondition holds}
\]
rather than that the procedure returns on every input or that a probabilistic
retry terminates almost surely. This distinction is essential because of
KeyGen and Sign rejection loops.

\subsection{How far algebraic geometry applies}

The ring \(R_q\) of course defines the affine scheme
\(\operatorname{Spec}R_q\). Local memory restrictions, local equalities,
and frame conditions also suggest restriction maps and gluing.
The current EasyCrypt files, however, do not formalize schemes, varieties,
spectra, sheaves, or cohomology. Proofs chiefly concern arrays over finite sets,
fixed-width words, state transitions, and integer division.

\begin{analogy}
Viewing memory \(m\) as a section over an address set, and equality on a region
\(U\) as \(m_1\restrict_U=m_2\restrict_U\), quickly suggests the shape of
region-local hash theorems. The analogy does not use sheaf axioms or a scheme
structure to derive EasyCrypt conclusions. The actual basis is pointwise
\identifier{load8} equality and the procedure's read footprint.
\end{analogy}

Algebraic geometers should pay particular attention to hashes.
SHAKE state transitions are generally neither \(R_q\)-linear maps nor regular
morphisms. This guide treats hashing as a deterministic state machine reading
finite byte strings. Equality of hash inputs therefore yields equality of
outputs through relational execution theorems about the two actual programs,
rather than through algebraic functionality.

\subsection{Six statements to remember first}

\begin{enumerate}
  \item Top-level \identifier{FC_KeyGen}, \identifier{FC_Sign}, and
        \identifier{FC_Verify} are still goal interfaces, not completed theorems.
  \item \(\prefix_{992}(sk)=vk\) and raw/internal \(\mu\)-prefix agreement
        are proved down to actual generated procedures, but do not automatically
        establish complete API security.
  \item The pure HBZ/rANS inverse structure, individual actual encoder/decoder
        trace refinements, core-harness inverse including actual copying,
        and production full-HBZ wrapper arrows are closed as
        success-conditioned partial correctness.
  \item For the all-6 stream generated by canonical all-zero HBZ, failure of a
        returning actual encoder is excluded and the production zero round
        trip is proved; fixed-input termination, losslessness, and
        \identifier{phoare} remain open.
  \item The actual KeyGen matrix/finalize snapshot has a proved KG-2-like
        decomposition and snapshot-only mod-\(2q\) congruence. At the frozen-paper
        stage, however, full-NTT convolution and the representation adapter
        following the final \identifier{output_row} NTT-word rewrite were
        absent, so work was frozen under \identifier{STOP-KG-NTT}.
        Sign was subsequently frozen under \identifier{STOP-SIGN-CHAL-MODE2};
        Verify was frozen at \FrozenPaperGate, lacking the actual 2-by-4
        NTT/CRT/freeze procedure
        \cite{week16-sign-report,week16-verify-report}.
        Post-freeze work closed pure convolution, the native \(R_q\) row
        product, and \(R_q\)--HAETAE representation transport, and interpreted
        actual \identifier{avec} as the paper's \(a\).
        But security-side synthetic ExpandVecA differs from actual
        SHAKE128/rejection semantics and the paper's \(qj\), so the current
        status is \CurrentGate. This does not change the frozen 82-target
        verdict \cite{post-freeze-kg-rq-haetae-report,
        post-freeze-kg-actual-avec-qj-report,paper-freeze-artifacts}.
  \item Separate final cryptographic assumptions from implementation proof
        obligations. Implementation correctness beyond the MLWE/BST-MSIS/ROM
        interfaces is not hidden as an assumption.
\end{enumerate}

These six statements govern status judgments throughout the following chapters.
